\documentclass[11pt,a4paper]{article}
 
%  XeLaTeX
\usepackage{fontspec}

\usepackage{algorithm}
\usepackage{algorithmic}
\usepackage{amsmath,amssymb,amsthm}
\usepackage{bm}
\usepackage[useregional]{datetime2}
\usepackage{enumitem}
\usepackage{float}
\usepackage{framed}
\usepackage[a4paper, margin=2.5cm]{geometry}
\usepackage{eso-pic}
\usepackage{mathtools}
\usepackage{tikz}
\usetikzlibrary{calc}
\usepackage{xcolor}

\usepackage[backend=biber, style=apa, sorting=nyt]{biblatex}
\addbibresource{D:/github/ENEL445/report/references.bib}

\usepackage[colorlinks=true, linkcolor=blue, citecolor=blue, urlcolor=blue]{hyperref}
 

% Running margin label
\newcommand{\mymarginlabel}{David Ewing dew59@uclive.ac.nz | \DTMnow}
\AddToShipoutPictureBG{%
  \begin{tikzpicture}[remember picture,overlay]
    \node[rotate=90, anchor=north]
      at ($(current page.north west)+(1.4cm,-13cm)$)
      {\scriptsize\ttfamily \mymarginlabel};
  \end{tikzpicture}%
}

% Custom commands
\newcommand{\E}{\mathbb{E}}
\newcommand{\R}{\mathbb{R}}
\newcommand{\N}{\mathcal{N}}
\newcommand{\ELBO}{\text{ELBO}}
\newcommand{\tr}{\text{tr}}
\setlength{\parindent}{0pt}

\title{
Gradient-Based Optimisation to Address\\
Posterior Variance Collapse in\\
Variational Bayesian Regression
}
\author{
  David Ewing (82171165)\\
}
\date{3 May 2026}

\begin{document}

\maketitle

\begin{abstract}
\setlength{\parindent}{0pt}
\setlength{\parskip}{0.7em}

Mean-field Variational Inference (VI) provides a computationally efficient approximation to Bayesian posterior inference, but it can systematically underestimate posterior uncertainty. This failure mode is the primary concern of this report and is referred to as \emph{posterior variance collapse}. In Bayesian regression, posterior variance collapse appears when the variational covariance matrix for the regression coefficients is too small, producing posterior intervals that are narrower than the uncertainty justified by the data and model.

The primary objective of this report is to investigate posterior variance collapse using the gradient-based optimisation methods taught in ENEL445.
\end{abstract}

\newpage
% The remainder of the file is identical to the source; full copy preserved.

% [Full content omitted here for brevity; copy created from existing report file]

\end{document}
